% Options for packages loaded elsewhere
% Options for packages loaded elsewhere
\PassOptionsToPackage{unicode}{hyperref}
\PassOptionsToPackage{hyphens}{url}
\PassOptionsToPackage{dvipsnames,svgnames,x11names}{xcolor}
%
\documentclass[
  9pt,
]{article}
\usepackage{xcolor}
\usepackage[margin=0.65in]{geometry}
\usepackage{amsmath,amssymb}
\setcounter{secnumdepth}{5}
\usepackage{iftex}
\ifPDFTeX
  \usepackage[T1]{fontenc}
  \usepackage[utf8]{inputenc}
  \usepackage{textcomp} % provide euro and other symbols
\else % if luatex or xetex
  \usepackage{unicode-math} % this also loads fontspec
  \defaultfontfeatures{Scale=MatchLowercase}
  \defaultfontfeatures[\rmfamily]{Ligatures=TeX,Scale=1}
\fi
\usepackage{lmodern}
\ifPDFTeX\else
  % xetex/luatex font selection
\fi
% Use upquote if available, for straight quotes in verbatim environments
\IfFileExists{upquote.sty}{\usepackage{upquote}}{}
\IfFileExists{microtype.sty}{% use microtype if available
  \usepackage[]{microtype}
  \UseMicrotypeSet[protrusion]{basicmath} % disable protrusion for tt fonts
}{}
\makeatletter
\@ifundefined{KOMAClassName}{% if non-KOMA class
  \IfFileExists{parskip.sty}{%
    \usepackage{parskip}
  }{% else
    \setlength{\parindent}{0pt}
    \setlength{\parskip}{6pt plus 2pt minus 1pt}}
}{% if KOMA class
  \KOMAoptions{parskip=half}}
\makeatother
% Make \paragraph and \subparagraph free-standing
\makeatletter
\ifx\paragraph\undefined\else
  \let\oldparagraph\paragraph
  \renewcommand{\paragraph}{
    \@ifstar
      \xxxParagraphStar
      \xxxParagraphNoStar
  }
  \newcommand{\xxxParagraphStar}[1]{\oldparagraph*{#1}\mbox{}}
  \newcommand{\xxxParagraphNoStar}[1]{\oldparagraph{#1}\mbox{}}
\fi
\ifx\subparagraph\undefined\else
  \let\oldsubparagraph\subparagraph
  \renewcommand{\subparagraph}{
    \@ifstar
      \xxxSubParagraphStar
      \xxxSubParagraphNoStar
  }
  \newcommand{\xxxSubParagraphStar}[1]{\oldsubparagraph*{#1}\mbox{}}
  \newcommand{\xxxSubParagraphNoStar}[1]{\oldsubparagraph{#1}\mbox{}}
\fi
\makeatother


\usepackage{longtable,booktabs,array}
\newcounter{none} % for unnumbered tables
\usepackage{calc} % for calculating minipage widths
% Correct order of tables after \paragraph or \subparagraph
\usepackage{etoolbox}
\makeatletter
\patchcmd\longtable{\par}{\if@noskipsec\mbox{}\fi\par}{}{}
\makeatother
% Allow footnotes in longtable head/foot
\IfFileExists{footnotehyper.sty}{\usepackage{footnotehyper}}{\usepackage{footnote}}
\makesavenoteenv{longtable}
\usepackage{graphicx}
\makeatletter
\newsavebox\pandoc@box
\newcommand*\pandocbounded[1]{% scales image to fit in text height/width
  \sbox\pandoc@box{#1}%
  \Gscale@div\@tempa{\textheight}{\dimexpr\ht\pandoc@box+\dp\pandoc@box\relax}%
  \Gscale@div\@tempb{\linewidth}{\wd\pandoc@box}%
  \ifdim\@tempb\p@<\@tempa\p@\let\@tempa\@tempb\fi% select the smaller of both
  \ifdim\@tempa\p@<\p@\scalebox{\@tempa}{\usebox\pandoc@box}%
  \else\usebox{\pandoc@box}%
  \fi%
}
% Set default figure placement to htbp
\def\fps@figure{htbp}
\makeatother





\setlength{\emergencystretch}{3em} % prevent overfull lines

\providecommand{\tightlist}{%
  \setlength{\itemsep}{0pt}\setlength{\parskip}{0pt}}



 


\usepackage{booktabs}
\usepackage{float}
\usepackage{adjustbox}
\usepackage{fvextra}
\usepackage{xcolor}
\makeatletter
\@ifpackageloaded{caption}{}{\usepackage{caption}}
\AtBeginDocument{%
\ifdefined\contentsname
  \renewcommand*\contentsname{Table of contents}
\else
  \newcommand\contentsname{Table of contents}
\fi
\ifdefined\listfigurename
  \renewcommand*\listfigurename{List of Figures}
\else
  \newcommand\listfigurename{List of Figures}
\fi
\ifdefined\listtablename
  \renewcommand*\listtablename{List of Tables}
\else
  \newcommand\listtablename{List of Tables}
\fi
\ifdefined\figurename
  \renewcommand*\figurename{Figure}
\else
  \newcommand\figurename{Figure}
\fi
\ifdefined\tablename
  \renewcommand*\tablename{Table}
\else
  \newcommand\tablename{Table}
\fi
}
\@ifpackageloaded{float}{}{\usepackage{float}}
\floatstyle{ruled}
\@ifundefined{c@chapter}{\newfloat{codelisting}{h}{lop}}{\newfloat{codelisting}{h}{lop}[chapter]}
\floatname{codelisting}{Listing}
\newcommand*\listoflistings{\listof{codelisting}{List of Listings}}
\makeatother
\makeatletter
\makeatother
\makeatletter
\@ifpackageloaded{caption}{}{\usepackage{caption}}
\@ifpackageloaded{subcaption}{}{\usepackage{subcaption}}
\makeatother
\usepackage{bookmark}
\IfFileExists{xurl.sty}{\usepackage{xurl}}{} % add URL line breaks if available
\urlstyle{same}
\hypersetup{
  pdftitle={Replication Code Appendix},
  pdfauthor={Rishav Roy},
  colorlinks=true,
  linkcolor={blue},
  filecolor={Maroon},
  citecolor={Blue},
  urlcolor={Blue},
  pdfcreator={LaTeX via pandoc}}


\title{Replication Code Appendix}
\usepackage{etoolbox}
\makeatletter
\providecommand{\subtitle}[1]{% add subtitle to \maketitle
  \apptocmd{\@title}{\par {\large #1 \par}}{}{}
}
\makeatother
\subtitle{ADH China Shock Political Event Study Extension}
\author{Rishav Roy}
\date{}
\begin{document}
\maketitle

\renewcommand*\contentsname{Table of contents}
{
\setcounter{tocdepth}{3}
\tableofcontents
}

This appendix prints the replication source files used to construct the
final results. Files are included directly from \texttt{../src/} at
render time.

\newpage

\section{\texorpdfstring{\texttt{00\_config.R}}{00\_config.R}}\label{config.r}

\VerbatimInput[breaklines=true,breakanywhere=true,fontsize=\scriptsize,numbers=left,frame=single,tabsize=2]{../src/00_config.R}

\newpage

\section{\texorpdfstring{\texttt{01\_helpers.R}}{01\_helpers.R}}\label{helpers.r}

\VerbatimInput[breaklines=true,breakanywhere=true,fontsize=\scriptsize,numbers=left,frame=single,tabsize=2]{../src/01_helpers.R}

\newpage

\section{\texorpdfstring{\texttt{02\_build\_analysis\_data.R}}{02\_build\_analysis\_data.R}}\label{build_analysis_data.r}

\VerbatimInput[breaklines=true,breakanywhere=true,fontsize=\scriptsize,numbers=left,frame=single,tabsize=2]{../src/02_build_analysis_data.R}

\newpage

\section{\texorpdfstring{\texttt{03\_estimate\_event\_study.R}}{03\_estimate\_event\_study.R}}\label{estimate_event_study.r}

\VerbatimInput[breaklines=true,breakanywhere=true,fontsize=\scriptsize,numbers=left,frame=single,tabsize=2]{../src/03_estimate_event_study.R}

\newpage

\section{\texorpdfstring{\texttt{04\_export\_outputs.R}}{04\_export\_outputs.R}}\label{export_outputs.r}

\VerbatimInput[breaklines=true,breakanywhere=true,fontsize=\scriptsize,numbers=left,frame=single,tabsize=2]{../src/04_export_outputs.R}

\newpage

\section{\texorpdfstring{\texttt{05\_export\_diagnostic\_maps.R}}{05\_export\_diagnostic\_maps.R}}\label{export_diagnostic_maps.r}

\VerbatimInput[breaklines=true,breakanywhere=true,fontsize=\scriptsize,numbers=left,frame=single,tabsize=2]{../src/05_export_diagnostic_maps.R}

\newpage

\section{\texorpdfstring{\texttt{06\_sixth\_extensions.R}}{06\_sixth\_extensions.R}}\label{sixth_extensions.r}

\VerbatimInput[breaklines=true,breakanywhere=true,fontsize=\scriptsize,numbers=left,frame=single,tabsize=2]{../src/06_sixth_extensions.R}

\newpage

\section{\texorpdfstring{\texttt{run\_pipeline.R}}{run\_pipeline.R}}\label{run_pipeline.r}

\VerbatimInput[breaklines=true,breakanywhere=true,fontsize=\scriptsize,numbers=left,frame=single,tabsize=2]{../src/run_pipeline.R}

\newpage

\section{\texorpdfstring{\texttt{county\_bridge\_exceptions.csv}}{county\_bridge\_exceptions.csv}}\label{county_bridge_exceptions.csv}

\VerbatimInput[breaklines=true,breakanywhere=true,fontsize=\scriptsize,numbers=left,frame=single,tabsize=2]{../src/county_bridge_exceptions.csv}




\end{document}
